% Source-faithful assembly: Mumford, Abelian Varieties, Appendix I,
% Horizon transcription pages 0262--0271.

\section{Applications of Tate's theorem}
\label{mumford-appendix-I-applications-source}
\dcref{M:appI:I}

\begin{proof}[Continuation of Lemma 5 and Step VII]
Let $F(X)=a_0X^n+a_1X^{n-1}+\cdots+a_n$. The assertion is trivial when
$a_n=0$, since then $X$ is a factor. After substituting $a_nX$ for $X$ and
removing the common factor, assume $a_n=1$. For a finite set of primes $S$ put
$N=\prod_{p\in S}p$. Then
\[
 F(vN)=a_0v^nN^n+\cdots+a_{n-1}vN+1\equiv1\pmod N,
\]
so no prime in $S$ divides $F(vN)$. For $v$ large, $F(vN)\ne\pm1$ and has a
prime factor outside $S$. This proves Lemma 5 and supplies infinitely many
such primes.

For all $l\ne {\rm char}\,k_0$, the dimension of
${\rm End}_{\mathbb Q_l}(V_l)^{(G)}$ is the same. Assume all endomorphisms of
$X$ are defined over $k_0$. Frobenius $\pi$ is central in ${\rm End}^0X$;
$\mathbb Q_l[\pi]$ is semisimple and the image $\pi'$ on $V_l$ is semisimple.
Its characteristic polynomial $P(t)$ has integral coefficients independent of
$l$. The closed subgroups generated by $\pi^{n!}$ form a fundamental system of
neighborhoods of the identity, so the invariant endomorphisms are the
commutant of $\pi^{n!}$ for large $n$. If $\theta_1,\ldots,\theta_{2g}$ are
the roots of $P$, the roots of the characteristic polynomial of $\pi^{n!}$
are $\theta_i^{n!}$; for large $n$ their distinct values and multiplicities
stabilize. The assertion follows from the next lemma.
\source{mumford-abelian-varieties:page-0262}
\end{proof}
\begin{lemma}[Commutants of semisimple endomorphisms]
\label{mumford-lem-appI-commutant-dimension}
Let $A$ and $B$ be absolutely semisimple endomorphisms of finite-dimensional
$k$-vector spaces $V,W$, with
\[
 P_A=\prod_p p^{m(p)},\qquad P_B=\prod_p p^{n(p)}
\]
the decompositions into powers of distinct irreducible monic polynomials. Then
\[
 E=\{\varphi\in{\rm Hom}_k(V,W):\varphi A=B\varphi\}
\]
has dimension
\[
 r(P_A,P_B)=\sum_p m(p)n(p)\deg p,
\]
and this integer is unchanged by extension of $k$.
\source{mumford-abelian-varieties:page-0262}
\end{lemma}
\begin{proof}
Make $V,W$ into $k[X]$-modules using $A,B$, and write
$M_p=k[X]/(p(X))$. Semisimplicity gives
\[
 V\simeq\prod_pM_p^{m(p)},\qquad W\simeq\prod_pM_p^{n(p)}.
\]
The $M_p$ are non-isomorphic simple modules, while
$E={\rm Hom}_{k[X]}(V,W)$. Since
$\dim_k{\rm Hom}_{k[X]}(M_p,M_p)=\dim_kM_p=\deg p$, the formula follows;
the same module calculation commutes with base extension.

\end{proof}


\begin{theorem}[Tate, strengthened form]
\label{mumford-thm-tate-hom-finite-field-rational}
For abelian varieties $X,Y$ over a finite field $k_0$,
\[
 \mathbb Z_l\otimes_{\mathbb Z}{\rm Hom}_{k_0}(X,Y)
 \xrightarrow{\sim}{\rm Hom}_{\mathbb Z_l}(T_l(X),T_l(Y))^G.
\]
\source{mumford-abelian-varieties:page-0263}
\uses{mumford-thm-tate-hom-finite-field,mumford-lem-appI-commutant-dimension}
\end{theorem}
\begin{proof}
Take $G$-invariants on both sides of the isomorphism (2).

\end{proof}


\begin{theorem}[Tate's applications]
\label{mumford-thm-tate-applications}
Let $X,Y$ be abelian varieties over a finite field $k_0$, with Frobenius
characteristic polynomials $P_X,P_Y$. Then:
\begin{enumerate}
\item $\operatorname{rank}{\rm Hom}_{k_0}(X,Y)=r(P_X,P_Y)$.
\item The following are equivalent: (b1) $Y$ is $k_0$-isogenous to an abelian
subvariety of $X$ defined over $k_0$; (b2) $V_l(Y)$ is $G$-isomorphic to a
$G$-subspace of $V_l(X)$ for some $l$; (b3) $P_Y$ divides $P_X$.
\item The following are equivalent: (c1) $X,Y$ are $k_0$-isogenous; (c2)
$P_X=P_Y$; (c3) they have the same number of $k_1$-rational points for every
finite extension $k_1/k_0$.
\end{enumerate}
\source{mumford-abelian-varieties:page-0264}
\uses{mumford-thm-tate-hom-finite-field-rational,mumford-lem-appI-commutant-dimension}
\end{theorem}
\begin{proof}
Part (a) follows from the strengthened theorem and the commutant lemma, since
Frobenius topologically generates $G$. The implications
$(b_1)\Rightarrow(b_2)\Rightarrow(b_3)$ are immediate. If (b2) holds, choose
an injective $G$-homomorphism $\varphi:T_l(Y)\to T_l(X)$. Approximate it by
$T_l(\psi)$ for $\psi\in{\rm Hom}_{k_0}(X,Y)$, retaining injectivity. If
$\psi$ were not an isogeny, an abelian subvariety in its kernel would give a
nonzero submodule $T_l(Z)$ in the kernel, contradiction. Thus $\psi$ is an
isogeny, proving (b).

The equivalence $(c_1)\Leftrightarrow(c_2)$ is the equal-dimension case of
(b). If $\omega_i,\omega'_j$ are roots of $P_X,P_Y$, then
\[
 N_n=\prod_i(1-\omega_i^n),\qquad N'_n=\prod_j(1-\omega_j^{\prime n}).
\]
The implication from equal polynomials is clear. Conversely, expand
\[
 \prod_{i\in I}(1-\omega_i^n)=\sum_{S\subset I}(-1)^{|S|}\omega_S^n,
 \quad
 \prod_{j\in J}(1-\omega_j^{\prime n})=\sum_{T\subset J}(-1)^{|T|}\omega_T^{\prime n}.
\]
After multiplying by $t^n$ and summing,
\[
 \sum_{S\subset I}(-1)^{|S|}\frac1{1-t\omega_S}
 =\sum_{T\subset J}(-1)^{|T|}\frac1{1-t\omega'_T}.
\]
Since all roots have absolute value $q^{1/2}$, comparison of poles on
$|t|=q^{-1/2}$ gives a bijection of the roots and $P_X=P_Y$.

\end{proof}


\begin{proposition}[Descent of closed abelian subvarieties]
\label{mumford-prop-closed-subvariety-descent-finite-field}
If $Y$ is an abelian subvariety of an abelian variety $X$ over a finite field
$k_0$ and is $k_0$-closed, then $Y$ is an abelian variety defined over $k_0$.
Moreover, there is a $k_0$-defined abelian subvariety $Z$ with $Y+Z=X$ and
$Y\cap Z$ finite. Consequently every $k_0$-abelian variety is isogenous to a
product of $k_0$-simple abelian varieties, and if
$X\sim\prod_iX_i^{n_i}$ with pairwise non-isogenous simple $X_i$, then
\[
 {\rm End}_{k_0}^{0}X\simeq\prod_iM_{n_i}(D_i),
 \qquad D_i={\rm End}_{k_0}^{0}(X_i).
\]
\source{mumford-abelian-varieties:page-0265}
\end{proposition}
\begin{proof}
Give the corresponding closed subset $Y_0$ of $X_0$ its reduced structure.
The multiplication restricts set-theoretically to $Y_0$, and over a finite
(hence perfect) field the relevant function fields are separable, so the base
change and product are reduced; hence $Y$ descends as an abelian variety.

For an ample line bundle $L$ on $X$ defined over $k_0$, the connected component
of zero of
\[
 Z'=\{z\in X:T_z^*L\otimes L^{-1}\text{ is trivial}\}
\]
is a complement up to finite intersection. The group $Z'$ is stable under
Galois, and its component containing the rational point $0$ is stable as well,
 so it descends. Repeating the complement argument gives the decomposition and
 the displayed endomorphism-algebra structure.
\source{mumford-abelian-varieties:page-0266}

\end{proof}


\begin{theorem}[Structure of endomorphism algebras]
\label{mumford-thm-endomorphism-algebra-finite-field}
Let $X$ have dimension $g$ over a finite field $k_0$, with Frobenius $\pi$ and
characteristic polynomial $P$. Put $F=\mathbb Q[\pi]$ and
$E={\rm End}_{k_0}^{0}(X)$. Then:
\begin{enumerate}
\item $F$ is the center of the semisimple algebra $E$.
\item $E$ contains a semisimple maximal commutative $\mathbb Q$-subalgebra
of rank $2g$.
\item The following are equivalent: $[E:\mathbb Q]=2g$; $P$ has no multiple
root; $E=F$; $E$ is commutative.
\item The following are equivalent: $[E:\mathbb Q]=(2g)^2$; $P$ is a power
of a linear polynomial; $F=\mathbb Q$; $E=M_g(D_p)$ for the quaternion division
algebra $D_p$ split away from $p,\infty$; $X$ is isogenous to the $g$th power
of a supersingular curve with all endomorphisms defined over $k_0$.
\item $X$ is isogenous to a power of a $k_0$-simple variety iff $P$ is a power
of a $\mathbb Q$-irreducible polynomial. In that case $E$ is central simple over
$F$, split at finite primes not dividing $p$ and nonsplit at real primes.
\end{enumerate}
\source{mumford-abelian-varieties:page-0267}
\uses{mumford-thm-tate-applications,mumford-prop-closed-subvariety-descent-finite-field}
\end{theorem}
\begin{proof}
The main theorem gives $F_1=\mathbb Q_l\otimes F$ as the center of
$E_1=\mathbb Q_l\otimes E$, proving (a). Write
$E=\prod_iA_i$, with centers $K_i$, $[K_i:\mathbb Q]=a_i$, and
$[A_i:K_i]=b_i^2$. Choose maximal commutative semisimple $L_i\supset K_i$
with $[L_i:K_i]=b_i$. Then $L=\prod_iL_i$ is maximal commutative and
$[L:\mathbb Q]=\sum_i a_ib_i$.

For any $l$, write $K_i\otimes\mathbb Q_l=\prod_jK'_{ij}$ and decompose
$P=\prod_\nu P_\nu^{m_\nu}$ over $\mathbb Q_l$. As a module,
\[
 T_l(X)\simeq\prod_\nu S_\nu^{m_\nu},\qquad
 S_\nu=\mathbb Q_l[T]/(P_\nu),
\]
so the commutant is $\prod_\nu M_{m_\nu}(S_\nu)$. Comparison gives splitting at
every finite prime away from $p$, a partition into the factors of each $A_i$,
and $m_\nu=b_i$ on the corresponding part. Therefore
\[
 \sum_i a_ib_i=\sum_\nu m_\nu[S_\nu:\mathbb Q_l]=\deg P=2g.
\]
This proves (b), and commutativity is equivalent to all $m_\nu=1$, proving (c).

The equality $[E:\mathbb Q]=(2g)^2$ is equivalent to
$E\otimes\mathbb Q_l=M_{2g}(\mathbb Q_l)$, hence to $P$ being a power of a
linear polynomial and to $F=\mathbb Q$. Brauer invariants then give
$E=M_g(D_p)$; the structural decomposition and the supersingular criterion
give the equivalence with (d).

Let $Q$ be the product of the distinct irreducible factors of $P$. Since
$Q(\pi)=0$ and the factors divide the minimal polynomial, $Q$ is that minimal
polynomial. Thus $E$ is simple exactly when $F\simeq\mathbb Q[T]/(Q)$ is a
field, exactly when $P$ is a power of an irreducible polynomial. The splitting
 and real-prime assertions follow from the Brauer invariants; in the nonsquare
 case pass to the quadratic extension, where the center becomes $\mathbb Q$ and
 use (d).
\source{mumford-abelian-varieties:page-0268}

\end{proof}


\begin{corollary}[Elliptic curves]
\label{mumford-cor-elliptic-cm-isogenous}
Any two elliptic curves over finite fields with isomorphic algebras of complex
multiplications are isogenous over the algebraic closure. In particular, any
two supersingular elliptic curves are isogenous.
\source{mumford-abelian-varieties:page-0269}
\uses{mumford-thm-tate-applications,mumford-thm-endomorphism-algebra-finite-field}
\end{corollary}
\begin{proof}
For supersingular $X,Y$, choose a common finite field $k_0$ over which their
endomorphism algebras are quaternion algebras, so both Frobenius morphisms lie
in $\mathbb Q$ and have squares equal to $q=|k_0|$. Over the quadratic
extension their Frobenius actions agree, so Theorem 2 gives an isogeny.

More generally, if both endomorphism algebras are an imaginary quadratic field
$K$, choose a common field over which all endomorphisms and $p$-torsion points
are rational. The $p$-adic representation splits $p$ as
$\mathfrak p\mathfrak p'$ in $K$, and (after conjugating one identification)
both Frobenius elements are powers of $\mathfrak p$. Their norms are the same
power of $p$, so they differ by a unit, hence by a root of unity. Suitable
powers have the same minimal, and therefore characteristic, polynomial; Theorem
2 gives an isogeny after a finite extension.
\source{mumford-abelian-varieties:page-0270}
\source{mumford-abelian-varieties:page-0271}

\end{proof}
